(1) 在地面系顺风传播的波面速度为 $c+w$；声源相邻两次发射相隔周期 $1/f_s$，其位置前进 $v_s/f_s$。因此两个波面沿 $+x$ 的空间间距为 $\lambda_+=(c+w-v_s)/f_s$。该式必须使用介质相对地面的波面速度，且假设声源向右的亚声速条件 $v_s<c+w$ 使前向波面排列正常。
(2) 接收器与顺向波面的相对速度为 $c+w-v_o$，每秒扫过的波面数即 $f_+=(c+w-v_o)/\lambda_+=f_s(c+w-v_o)/(c+w-v_s)$。当 $v_s=v_o=0$，虽然有风，频率仍为 $f_s$；当源和接收器以同一地面速度同向运动时比值为一。顺风会改变量测的波长，却不会单凭均匀介质流自动改变静止源与静止接收器的频率。
(3) 逆风向左波面的地面速度为 $w-c=-(c-w)$。源相邻发射时向右移 $v_s/f_s$，左向波面间距为 $\lambda_-=(c-w+v_s)/f_s$；向右移动的接收器迎着波面，扫过速度为 $c-w+v_o$。所以 $f_-=f_s(c-w+v_o)/(c-w+v_s)$，并与第二问在变换传播方向后符号规则一致。
(4) 经典声学存在空气静止系，波速 $c$ 相对于该介质定义，因此地面风速与源、接收器速度以伽利略方式组合。真空光传播不存在对应的优选静止介质，需使用洛伦兹变换与相对论多普勒公式。为避免常规波面穿越或激波堆积，上述线性周期公式要求顺向 $v_s<c+w$，逆向 $v_s>-(c-w)$，接收器也应在有关波面同向穿越的正常测速范围。
